% ===================================================================
%  TABEL Nr. 4 — Küps iga ja vanadus.
%  Traduction 2026 d'apres texto/io, controlee sur texto/fr et
%  texto/fi. Consigne : voir texto/et/10-tabel-01.tex.
%
%  CETTE PAGE REPOND A DEUX TABLEAUX ROMANCHES A LA FOIS, AU 4 ET AU
%  7, ET ELLE LES REUNIT.
%
%  Le tableau 4 romanche avait trouve que les mots de parente s'y
%  partagent entre l'adresse et la mention : « bab », « mamma »,
%  « tat » sont de nourrice parce qu'on les CRIE ; « abiadi »,
%  « quinà », « niera » sont latins parce qu'on n'en appelle
%  personne. L'ESTONIEN NE CONNAIT PAS CE PARTAGE : ses mots de
%  reference sont aussi natifs que ses vocatifs. « äi » le
%  beau-pere, « ämm » la belle-mere, « minia » la bru, « väimees »
%  le gendre, « küdi » le frere du mari, « nadu » la soeur du mari
%  — pas un emprunt, alors que huit siecles d'allemand ont laisse un
%  tiers du reste du lexique.
%
%  ET LE TABLEAU 7 ROMANCHE DONNE LA RAISON. Il y avait montre que le
%  romanche a un mot plein pour chaque case du betail — « bov »,
%  « tor », « vatga », « vadè » — la ou l'ido derive d'un seul
%  radical, et que LE NOMBRE DE RACINES SEPAREES MESURE LA DISTANCE A
%  LAQUELLE ON SE TIENT DE LA CHOSE.
%
%  L'estonien fait la meme chose sur la parente par alliance. La ou
%  l'allemand compose (Schwiegertochter, Schwiegersohn) et ou le
%  francais couvre quatre relations d'un seul « beau-frere »,
%  l'estonien a un mot plein pour chacune. Dans une ferme ou trois
%  generations vivent sous un toit, la difference entre « minia » et
%  « nadu » est la difference entre deux personnes qui doivent
%  s'entendre tous les jours.
%
%  LA LOI EST DONC LA MEME DES DEUX COTES DE L'EUROPE, ET ELLE NE
%  REGARDE PAS LA FAMILLE DE LANGUES : ON A UNE RACINE PLEINE POUR CE
%  QU'ON DOIT DISTINGUER EN AGISSANT, ET UN COMPOSE POUR CE QU'ON NE
%  DISTINGUE QU'EN PARLANT. Le romanche l'a montre sur l'etable,
%  l'estonien le montre sur la maisonnee.
%
%  « bandoliera vildo-sako » de l'alinea 4 est mis TOUT ENTIER en
%  gras, comme dans les vingt-deux colonnes precedentes.
%
%  UNE INVERSION, AU BLOC t04-c1-18-1. L'ido range sommelier (50),
%  bouteille (52), tire-bouchon (51) : l'instrument vient apres
%  l'objet. La premiere redaction estonienne mettait l'instrument
%  d'abord, comme l'estonien le fait volontiers, et renvoji.py l'a
%  signalee. L'ordre de l'ido tient sans forcer : le complement
%  d'objet passe devant et le comitatif suit.
% ===================================================================

%%K t04-apar-1 apar
\VUsaut{4.0mm}
\VUtitre{15.0pt}{60}{TABEL Nr. 4}
\VUsaut{1.6mm}
\VUtitre{11.4pt}{0}{Küps iga ja vanadus.}
\VUsaut{3.2mm}

%%K t04-tit-1 sub
\VUcentre{11.4pt}{0}{\textit{Esimene stseen.}}
\VUsaut{1.6mm}
\VUtitre{11.4pt}{0}{Pulmasöök.}
\VUsaut{2.6mm}

%%K t04-tit-2 sub
\VUpk{10.2pt}{40}{Sügis.}
\VUsaut{3.0mm}

%%K t04-c1-01-1 p
1. --- Noored on suuremaks kasvanud. Üks neist, Jaan, abiellub. Me viibime
\VUgras{pulmasöögil}, mis kogub sama laua ümber pruutpaari perekonnad.

%%K t04-c1-02-1 p
2. --- Me oleme \VUgras{sügisel}, \VUgras{septembrikuus}. \VUgras{Oktoobri} ja
\VUgras{novembri} külm tuul ei ole veel maad tema rohelisest lehestikust paljaks
teinud. Õhtune õhk on leige ja lõhnav; seepärast toimub õhtusöök aias
\VUgras{veranda} \textsuperscript{(8)} all.

%%K t04-c1-03-1 p
3. --- Kaugemal, \VUgras{künkal} \textsuperscript{(3)}, asub Jaani vanemate kodu,
elegantne \VUgras{loss} \textsuperscript{(1)}, mis on rohelusse peidetud; ilusad
viinamäed, mis annavad suurepärast veini, katavad künka nõlva. Viinamäekasvataja
harib ja hooldab neid.

%%K t04-c1-04-1 p
4. --- Viinamägede kohal lendab \VUgras{nurmkanade} \textsuperscript{(2)} parv, keda on
hirmutanud \VUgras{jahivahi} \textsuperscript{(5)} lähenemine. See viimane, \VUgras{püss}
kaenla all ja \VUgras{üle õla riputatud jahikott} seljas, uurib \VUgras{võsa}
\textsuperscript{(4)} koos oma \VUgras{koeraga} \textsuperscript{(6)}. Ta valvab mõisa ja
takistab salaküttidel ulukeid hävitamast.

%%K t04-c1-05-1 p
5. --- Väljaspool \VUgras{verandat} ronib \VUgras{viinapuu-spalee}, millelt ripuvad
tihedad \VUgras{viinamarjakobarad} \textsuperscript{(7)}.

%%K t04-c1-06-1 p
6. --- Ilusad sügislilled, mis kaunistavad \VUgras{lauda} \textsuperscript{(16)}, on
\VUgras{krüsanteemid} \textsuperscript{(9)}; pruudi \VUgras{isa} \textsuperscript{(10)} on
ühe neist pannud oma fraki nööpauku.

%%K t04-c1-07-1 p
7. --- Söök hakkab lõppema. \VUgras{Teener} \textsuperscript{(11)} hakkab tooma
magustoidutaldrikuid ja võtab peagi ära lauanõude eri osad, \VUgras{lusikad}
\textsuperscript{(14)}, \VUgras{kahvlid} \textsuperscript{(12)}, \VUgras{noad}
\textsuperscript{(13)}, \VUgras{vaagnad} \textsuperscript{(15)}, mida enam ei vajata.

%%K t04-tit-3 sub
\VUpk{10.2pt}{40}{Sugulased.}
\VUsaut{3.0mm}

%%K t04-c1-08-1 p
8. --- Kõik näivad rõõmsad, ja naeratus, millega peigmehe \VUgras{ema}
\textsuperscript{(17)} oma lapsi vaatab, tõendab tema õnne.

%%K t04-c1-09-1 p
9. --- See abielu ühendab kaks maa rikast perekonda. Jaan, \VUgras{peigmees}
\textsuperscript{(18)}, on suure maaomaniku \VUgras{poeg} ja temast saab osava
tööstureni \VUgras{väimees}.

%%K t04-c1-10-1 p
10. --- \VUgras{Pruutpaar} on noor ja siiski olid nad ammu \VUgras{kihlatud}.
\VUgras{Noorik} \textsuperscript{(19)}, Martha, on võluv oma \VUgras{valges kleidis},
mida kaunistavad apelsinipuu õied.

%%K t04-c1-11-1 p
11. --- Tema \VUgras{äi} \textsuperscript{(20)} on väga õnnelik, et tal on nii lahke
\VUgras{minia}, ja Jaani \VUgras{ämm} \textsuperscript{(21)} naeratab, nähes oma
\VUgras{tütart} nii ilusana.

%%K t04-c1-12-1 p
12. --- Laua otsas istuvad ülejäänud \VUgras{külalised} \textsuperscript{(22)},
\VUgras{sugulased, sõbrad, onupojad, onutütred}. \VUgras{Peakelner}
\textsuperscript{(23)}, salvrätik käe all, teenindab neid usinasti.

%%K t04-tit-4 sub
\VUpk{10.2pt}{40}{Sõbrad.}
\VUsaut{3.0mm}

%%K t04-c1-13-1 p
13. --- Laua paremal pool on \VUgras{preili} \textsuperscript{(24)}, noore pruudi
\VUgras{sõbratar}, siis tema \VUgras{vend}, kes on tema \VUgras{peiupoiss}
\textsuperscript{(25)}. Ta vestleb oma \VUgras{pruutneitsiga} \textsuperscript{(26)}, kes
on peigmehe õde ja kellest saab selle abieluga Martha \VUgras{nadu}. Ta on võluv noor
neiu. Tal on ilusad \VUgras{ehted}, \VUgras{pärlikee} ja kuldne \VUgras{käevõru}.

%%K t04-c1-14-1 p
14. --- Nende lähedal on härra, kes seisab ja tõstab oma klaasi, perekonna
\VUgras{sõber} \textsuperscript{(27)}. Ta on üks \VUgras{peigmehe tunnistajaist}. Ta
\VUgras{peab toosti}, joob noore pruutpaari \VUgras{terviseks} ja soovib neile pikka
õnne. Ta on riietatud pidulikku ülikonda, frakki ja \VUgras{lakknahast kingadesse}
\textsuperscript{(28)}. Ta saadab \VUgras{vanaema} \textsuperscript{(29)}, kes on
\VUgras{lesk}.

%%K t04-c1-15-1 p
15. --- Noormees \VUgras{frakis} \textsuperscript{(31)}, peene musta vuntsiga, on Jaani
\VUgras{küdi} \textsuperscript{(30)}. Ta on väljapaistev insener. Ta on väga elegantse
\VUgras{noore daami} \textsuperscript{(33)} naaber, kes on Martha \VUgras{vanem õde} ja
selle armsa tüdrukukese ema, kes tuleb, andes käe oma onupojale, pakkuma lille ja
\VUgras{soovima} talle \VUgras{head õhtut}. Sellel daamil on väga ilusad
\VUgras{kõrvarõngad} ja elegantne \VUgras{peakate}.

%%K t04-c1-16-1 p
16. --- Väike onupoeg, kes on oma \VUgras{pluusi} \textsuperscript{(36)} ja oma pungis
\VUgras{põlvpükste} \textsuperscript{(37)} pärast nii kummaline, on väga haletsusväärne.
Ta on \VUgras{orb} \textsuperscript{(35)}. Väga noorena kaotas ta oma isa ja oma ema.
Tema onu on tema \VUgras{eestkostja} \textsuperscript{(38)} ja jäi vallaliseks, et
vaest last kasvatada. Ta on heasüdamlik mees. Ta on pisut kiilas ja väga lühinägelik;
ta kasutab \VUgras{näpitsprille}.

%%K t04-tit-5 sub
\VUsaut{2.4mm}
\VUpk{10.2pt}{40}{Magustoit.}
\VUsaut{3.0mm}

%%K t04-c1-17-1 p
17. --- Pidulistest tagapool, \VUgras{laudlinaga} kaetud laual, on valmis pandud
\VUgras{magustoit} \textsuperscript{(39)}. Suures vaasis on siin puuvilju,
\VUgras{virsikud} \textsuperscript{(40)}, \VUgras{pirnid} \textsuperscript{(41)},
\VUgras{õunad} \textsuperscript{(42)}, \VUgras{aprikoosid} \textsuperscript{(43)} ja
\VUgras{viinamarjad} \textsuperscript{(44)}. Lähedal on \VUgras{ananassid}
\textsuperscript{(45)}, ilus \VUgras{kook} \textsuperscript{(46)},
\VUgras{likööripudelid} \textsuperscript{(48)} ja \VUgras{šampanja} \textsuperscript{(49)}
ning ka puhaste \VUgras{taldrikute} \textsuperscript{(47)} sambad.

%%K t04-c1-18-1 p
18. --- \VUgras{Keldrimees} \textsuperscript{(50)} avab vana veini \VUgras{pudeli}
\textsuperscript{(52)} \VUgras{korgitseriga} \textsuperscript{(51)}.
\VUgras{Korkkork} \textsuperscript{(53)} näib väga raskesti väljatõmmatav. Sellel
keldrimehel on \VUgras{salvrätik} \textsuperscript{(54)} käe all.

%%K t04-c1-19-1 p
19. --- Pärast õhtusööki lähevad härrad suitsetamistuppa ja daamid lähevad end balliks
valmis seadma.

%%K t04-tit-6 sub
\VUsaut{5.0mm}
\VUcentre{11.4pt}{0}{\textit{Teine stseen.}}
\VUsaut{1.6mm}
\VUpk{11.4pt}{40}{Vanaisa sünnipäev.}
\VUsaut{2.6mm}

%%K t04-tit-7 sub
\VUpk{10.2pt}{40}{Külaskäik.}
\VUsaut{3.0mm}

%%K t04-c2-01-1 p
1. --- Kümme aastat on möödunud, Jaani vanemad on vanaks jäänud ja armastavad väga oma
kolme lapselast, kaht \VUgras{pojapoega} \textsuperscript{(2)} ja üht
\VUgras{tütretütart} \textsuperscript{(3)}. Kuna täna on \VUgras{vanaisa sünnipäev},
tulevad lapsed neile head sünnipäeva soovima.

%%K t04-c2-02-1 p
2. --- Väike Peeter on veel väga noor, ta on ainult kolmeaastane. Ta näeb \VUgras{kassi}
\textsuperscript{(1)} lähenemas. Ta tahab silitada tema ilusat siidist \VUgras{karva} ja
tema pikka \VUgras{saba}; ta ei mõtle sellele, et graatsiliste \VUgras{käppade} all
peituvad kurjad teravad \VUgras{küüned}, mis teda ehk kriimustavad.

%%K t04-c2-03-1 p
3. --- Tema õde Gabriela, kes annab talle käe, on palju tõsisem, ja Marcel oma suure
\VUgras{mantli} \textsuperscript{(4)} ja nahkse \VUgras{vööga} \textsuperscript{(5)} näib
pisut mehistunud, vaatamata oma lühikestele pükstele, mis lasevad näha tema
\VUgras{sukki} \textsuperscript{(6)}.

%%K t04-c2-04-1 p
4. --- Nende isa on selga pannud oma paksu talvise \VUgras{palitu}
\textsuperscript{(7)} \VUgras{karusnahkse kraega} \textsuperscript{(8)} ja tema
\VUgras{taskus} \textsuperscript{(9)} on rätik. Tema abikaasal on viltkübar
\VUgras{sulgedega} \textsuperscript{(10)} kaunistatud, tema \VUgras{jakk}
\textsuperscript{(11)}, tema \VUgras{boa} \textsuperscript{(12)} ja tema \VUgras{muhv}
\textsuperscript{(13)}.

%%K t04-c2-05-1 p
5. --- On väga külm, sest valitseb talv. \VUgras{Lumi} \textsuperscript{(19)} on sadanud
terve päeva ja majade katused on täiesti valged. Koos \VUgras{ööga} muutub külm veel
teravamaks. Taevas säravad \VUgras{kuu} \textsuperscript{(17)} ja \VUgras{tähed}
\textsuperscript{(18)} ja läbistavad \VUgras{pimeduse}.

%%K t04-tit-8 sub
\VUsaut{4.2mm}
\VUpk{10.2pt}{40}{Haigus.}
\VUsaut{3.0mm}

%%K t04-c2-06-1 p
6. --- Vaene \VUgras{pime} \textsuperscript{(20)}, kes ületab platsi \VUgras{apteegi}
\textsuperscript{(16)} ees, juhitud oma \VUgras{puudlist} \textsuperscript{(21)}, on väga
õnnetu. Justiina, \VUgras{teenija} \textsuperscript{(14)}, toob köögist \VUgras{kausi}
\textsuperscript{(15)} väga kuuma \VUgras{taimeteed}, mida on heldelt suhkrustatud.

%%K t04-c2-07-1 p
7. --- \VUgras{Vanaema} \textsuperscript{(23)}, kes tahab otsida laualt oma
\VUgras{näpitsprille} \textsuperscript{(26)}, et lugeda \VUgras{retsepti}
\textsuperscript{(27)}, mille \VUgras{arst} \textsuperscript{(25)} on äsja kirjutanud,
tõuseb oma tugitoolist, toetudes oma \VUgras{kepile} \textsuperscript{(24)}.

%%K t04-c2-08-1 p
8. --- Sageli kannatab ta \VUgras{tõbede, peapöörituse, nohude, hambavalude} käes, tema
jalad on nõrgad ja tema silmad ei ole enam väga head. Sellest hoolimata pilkab ta
meelsasti oma vana \VUgras{arsti}, kes tahab talle alati \VUgras{rohusegu}
\textsuperscript{(28)} määrata. Täna on ta väga rõõmus, sest tema noor perekond on tema
ümber.

%%K t04-c2-09-1 p
9. --- Hästi suletud toas on mugav. Marmorist \VUgras{kaminas} \textsuperscript{(29)}
leegitseb \VUgras{kaunis tuli} \textsuperscript{(30)}, ja end tuntakse õnnelikuna
\VUgras{lambi} \textsuperscript{(31)} pehmes \VUgras{valguses}, mille on äsja süüdatud.

%%K t04-tit-9 sub
\VUsaut{4.2mm}
\VUpk{10.2pt}{40}{Vanaisa.}
\VUsaut{3.0mm}

%%K t04-c2-10-1 p
10. --- Isegi \VUgras{vanaisa} \textsuperscript{(33)} on unustanud, et ta oli haige, et
tema \VUgras{reuma} kinnitab ta oma \VUgras{tugitooli} \textsuperscript{(34)} külge ja et
alles eile oli tal \VUgras{palavik ja minestused}. Ta ei mäleta enam, et ta möödunud
talvel põdes \VUgras{kopsupõletikku} ja et kähe ja piinav \VUgras{köha} pani ta väga
kannatama.

%%K t04-c2-10-2 p
Ja siiski sai ta tervise tagasi ainult väga vaevaliselt. Nüüd on tal pisut parem. Kuid
tema jalg, mida ta toetab \VUgras{padjale} \textsuperscript{(38)}, mis on pandud väiksele
\VUgras{pingile} \textsuperscript{(39)}, teeb talle samuti valu.

%%K t04-c2-11-1 p
11. --- Kui ta on hoolikalt mähitud oma sooja \VUgras{hommikumantlisse}
\textsuperscript{(35)} paksude pärlmutrist \VUgras{nööpidega} \textsuperscript{(36)} ja
kui tal on enda kõrval, et talle ajalehte lugeda, väike \VUgras{orvuke}
\textsuperscript{(40)}, riietatud valgesse \VUgras{tanusse}, sinisesse \VUgras{kleiti}
\textsuperscript{(42)} ja \VUgras{pelerini} \textsuperscript{(41)}, siis ta ei kurda.

%%K t04-c2-12-1 p
12. --- Igal aastal, kui \VUgras{kalender} \textsuperscript{(43)} taas näitab esimese
\VUgras{detsembri} \VUgras{kuupäeva}, ootab ta kannatamatult oma \VUgras{lapselapsi},
sest ta teab, et nad ei unusta seda tähtsat \VUgras{kuupäeva}.

%%K t04-c2-13-1 p
13. --- Ta meenutab liigutatult väga ammust aega, mil ta ise läks head püha soovima
kuulsale keiserlikule \VUgras{marssalile}, kelle \VUgras{portree} \textsuperscript{(44)}
ripub seinal ja kes on tema laste \VUgras{vanavanaisa}.

\VUsaut{6.0mm}
\VUfilet{22mm}
